\documentclass[11pt]{article}
\usepackage[margin=1.05in]{geometry}
\usepackage{amsmath,amssymb,amsthm}
\usepackage{enumitem}
\usepackage[colorlinks=true,linkcolor=blue!50!black,citecolor=blue!50!black]{hyperref}
\usepackage{booktabs}

\newcommand{\R}{\mathbb{R}}
\newcommand{\E}{\mathbb{E}}
\newcommand{\Prob}{\mathbb{P}}
\newcommand{\dd}{\,\mathrm{d}}
\newcommand{\cH}{\mathcal{H}}
\newcommand{\cF}{\mathcal{F}}
\newcommand{\cR}{\mathcal{R}}
\newcommand{\one}{\mathbf{1}}
\newcommand{\code}[1]{\texttt{#1}}
\newtheorem{assumption}{Assumption}
\newtheorem{proposition}{Proposition}
\newtheorem{remark}{Remark}
\numberwithin{assumption}{section}
\numberwithin{proposition}{section}
\numberwithin{remark}{section}

\title{\bfseries The Escalation Ladder for Funding Events:\\[3pt]
A Mathematical Treatment}
\author{Counting processes, survival analysis, censoring, and self-excitation\\
for the staged modeling strategy of \code{STRATEGY.md}}
\date{}

\begin{document}
\maketitle

\begin{abstract}\noindent
This document gives the mathematical foundations for the model escalation
ladder: (1) inhomogeneous Poisson intensities driven by covariates; (2)
multi-state models with constrained stage transitions; (3) history-dependent
intensities through predictable covariates (``Hawkes without Hawkes''); (4)
fully self-exciting point processes and their cluster/cascade representation;
(5) the two-step treatment of firms outside the observed universe. For each
step we state the model in the counting-process language of survival analysis,
derive the likelihood and its properties, treat right censoring, interval
censoring and left truncation explicitly, list the assumptions in falsifiable
form, state what the data must exhibit for calibration to be meaningful, and
catalogue the statistical failure modes one should expect to find while
calibrating. The status-quo gradient-boosted benchmark is deliberately not
treated: it is a scoring rule, not an intensity model, and the mathematics
below is precisely what it lacks.
\end{abstract}

\tableofcontents

% =======================================================================
\section{Preliminaries: counting processes, intensities, and censoring}
\label{sec:prelim}

\subsection{The canonical setup}
Let $(\Omega,\cF,\Prob)$ carry a filtration $\{\cF_t\}$ representing
everything observed up to time $t$: past deals, covariate paths, and entries
and exits from the panel. For each firm $i$ let $N_i(t)$ count its funding
events in $[0,t]$. We work throughout with the \emph{conditional intensity}
\begin{equation}
\lambda_i(t)\dd t \;=\; \Prob\big(N_i \text{ jumps in } [t,t+\dd t) \mid \cF_{t-}\big),
\label{eq:intensity_def}
\end{equation}
equivalently the predictable process such that
$M_i(t)=N_i(t)-\int_0^t \lambda_i(s)\dd s$ is a local martingale
(Doob--Meyer). All estimators in this document are martingale-based in the
sense of Aalen's multiplicative intensity model \citep{aalen1978},
systematised in \citet{abgk1993}: we write
\begin{equation}
\lambda_i(t) \;=\; Y_i(t)\,\alpha_i(t),
\label{eq:aalen}
\end{equation}
where $Y_i(t)\in\{0,1\}$ is the \emph{at-risk indicator} (predictable, i.e.\
determined by $\cF_{t-}$) and $\alpha_i(t)$ the hazard of interest. The
at-risk process is where every observation problem of this dataset lives:
who is in the market, at which stage, tracked or not.

\paragraph{Point-process likelihood.} If a process with intensity
$\lambda_\vartheta$ is observed on $[0,T]$ with event times
$t_1<\dots<t_n$, the log-likelihood is
\begin{equation}
\ell(\vartheta)\;=\;\sum_{k=1}^{n}\log\lambda_\vartheta(t_k)
\;-\;\int_0^T \lambda_\vartheta(s)\dd s ,
\label{eq:pp_lik}
\end{equation}
\citep{daleyvj2003,ogata1978}. Everything on the ladder is an instance of
\eqref{eq:pp_lik} with a different parametrization of $\lambda$, estimated
either from exact times or after aggregation.

\subsection{A censoring taxonomy for funding data}
\label{sec:censoring}
The word ``censoring'' covers five distinct phenomena here; conflating them
is the single most common source of silent bias. For each we give the
likelihood treatment and the assumption under which that treatment is valid.

\begin{enumerate}[leftmargin=1.6em]
\item \textbf{Administrative right censoring.} Observation stops at the
sample end $T$; firms still active simply contribute the survival factor
$\exp(-\int_{t_{\text{last}}}^{T}\lambda_i)$ through \eqref{eq:pp_lik}. This
is Type-I censoring; it is \emph{always} independent and needs no assumption
beyond correct $\lambda$.
\item \textbf{Random right censoring by observed exit.} A firm leaves the
panel at an observed time $C_i$ (known IPO/M\&A/shutdown): set $Y_i(t)=0$
for $t>C_i$. Valid iff censoring is \emph{independent}: conditional on
covariates, the exit time carries no information about the counterfactual
funding hazard,
\begin{equation}
\lambda_i\big(t \mid \cF_{t-},\, C_i \ge t\big)
= \lambda_i\big(t \mid \cF_{t-}\big).
\label{eq:indep_censoring}
\end{equation}
For funding data this is \emph{not} innocuous --- firms exit \emph{because}
of funding states (graduation after success, death after drought) --- but it
is testable in the direction that matters: including exit-predictive
covariates in $\lambda$ restores validity (censoring independent
\emph{given covariates}, \citealp{kalbfleisch2002}).
\item \textbf{Informative, \emph{unobserved} right censoring.} The dangerous
one: a firm dies or graduates \emph{silently}, so $Y_i(t)$ is recorded as
$1$ while the truth is $0$. The at-risk process is then misspecified, the
denominator of every hazard is inflated, and --- crucially --- the bias is
\emph{directional}: covariates correlated with silent exit (staleness, late
stage, large cumulative raise) absorb the missing exits and can flip sign.
No reweighting of observables repairs this in general, because the
observation process is confounded with the outcome; the honest treatments
are (a) explicit exit modeling as a competing risk with a cure/mover--stayer
component \citep{farewell1982}, (b) bounds: report under the naive and the
best exit-labeled risk set, and quote the spread. This phenomenon is the
mathematical content of our measured ``risk-set inversion.''
\item \textbf{Interval censoring.} If only bucketed counts are kept (weekly
aggregation), event times are interval-censored. For a Poisson-type model
with intensity constant within buckets, the bucket counts are
\emph{sufficient}: the likelihood \eqref{eq:pp_lik} collapses exactly to a
Poisson GLM on counts (\S\ref{sec:step1}). For history-dependent models the
collapse is not exact, and what is lost is quantifiable Fisher information
about timescales (\S\ref{sec:step4}); for the full Hawkes case the
counts-only route is the interval-censored theory of \citet{rizoiu2022}.
\item \textbf{Left truncation (delayed entry).} A firm enters the observable
universe only at its first recorded deal (or first tracking date). Its
history before entry is unobservable, and it must not contribute risk before
entry: $Y_i(t)=0$ for $t<E_i$. Standard delayed-entry survival methodology
\citep{kalbfleisch2002} handles this \emph{if entry is independent given
covariates}; the deeper problem is that \emph{covariates themselves} are
undefined pre-entry --- the mathematical root of the newcomer channel
(\S\ref{sec:step5}), where the entrant's unobserved covariates are
integrated out rather than imputed.
\end{enumerate}

\begin{remark}[Coarsening]
Undisclosed deal sizes and estimated fields are a sixth, milder phenomenon:
\emph{coarsened} marks. Likelihood factorizations survive under Coarsening
At Random \citep{heitjan1991}; a disclosure indicator as a covariate is the
cheap guard when CAR is doubtful.
\end{remark}

% =======================================================================
\section{Step 1: covariate-driven Poisson intensities}
\label{sec:step1}

\subsection{Model and likelihood}
The first calibrated rung is the inhomogeneous Poisson (Cox-regression-style
multiplicative) model at whatever aggregation level is chosen --- sector
$s$, or firm $i$ within \eqref{eq:aalen}:
\begin{equation}
\lambda_s(t)\;=\;\exp\!\big(a_s+\beta_s^{\top}X(t)\big),
\label{eq:nhpp}
\end{equation}
with $X(t)$ exogenous (macro) covariates. Substituting into
\eqref{eq:pp_lik}, $\ell$ is concave in $(a_s,\beta_s)$ because the
log-intensity is affine and $-\int e^{\eta}$ is concave in $\eta$
(exponential-family argument; conditions and boundary cases in
\code{MATH\_NOTES.md} M2). MLE consistency and asymptotic normality for
point-process likelihoods of this type are classical \citep{ogata1978}.

\paragraph{Exact reduction under interval censoring.} If $X$ is constant on
buckets $[\tau_{w-1},\tau_w)$ (as-of weekly covariates), then
\eqref{eq:pp_lik} depends on the data only through bucket counts
$Y_{s,w}$, and maximizing it is \emph{identical} to the Poisson GLM
\begin{equation}
Y_{s,w}\sim\mathrm{Poisson}\!\big(\Delta_w\,e^{a_s+\beta_s^\top X_w}\big).
\label{eq:pois_glm}
\end{equation}
This equivalence is why step 1 loses nothing to weekly data: for a
history-free intensity, counts are sufficient statistics. It is the
\emph{only} rung with that property.

\subsection{Censored or not?}
Under \eqref{eq:nhpp} at sector level there is no firm at-risk process, so
censoring reduces to items 1 and 4 of \S\ref{sec:censoring}: administrative
truncation at $T$ (handled by \eqref{eq:pp_lik} automatically) and interval
censoring (handled exactly by \eqref{eq:pois_glm}). At the \emph{firm} level
the same model already inherits the full at-risk machinery, and every
warning of \S\ref{sec:censoring} item 3 applies --- which is one reason to
begin at sector level, where the aggregate is insensitive to which
individual firms are silently dead.

\subsection{Assumptions}
\begin{assumption}[No history dependence]\label{as:nhpp_indep}
Conditional on $\{X(s)\}_{s\le t}$, increments of $N$ are independent of
$\cF_{t-}$: $\lambda(t\mid\cF_{t-})=\lambda(t\mid X(t))$.
\end{assumption}
\begin{assumption}[Covariate exogeneity]\label{as:nhpp_exo}
$X$ is $\cF_{t-}$-predictable and its law does not depend on the parameters
of $N$ (no feedback from deals to macro at the horizon modeled).
\end{assumption}
\begin{assumption}[Correct link and log-linearity]\label{as:nhpp_link}
$\log\lambda$ is affine in $X$; departures are absorbed into diagnostics,
not the estimator.
\end{assumption}

\subsection{What the data must show, and where calibration fails}
\emph{To see in data:} quasi-Poisson dispersion $\approx 1$ per cell;
uniform randomized PIT of held-out counts; martingale residuals
$\widehat M_s(w)$ without autocorrelation. Each violation is informative:
overdispersion + PIT autocorrelation is the signature that
Assumption~\ref{as:nhpp_indep} fails and history matters (the gate to
steps 2--3); overdispersion \emph{without} autocorrelation indicates static
heterogeneity, pointing to a gamma-frailty/negative-binomial correction
\citep{vaupel1979} rather than to dynamics.

\emph{Statistical failure modes while calibrating:}
(i) \emph{collinearity} of macro series (rates, inflation, unemployment
comove) inflating $\mathrm{Var}(\hat\beta)$ --- ridge-penalize and report
the effective dof;
(ii) \emph{separation} in low-volume sectors (all-zero cells drive
$a_s\to-\infty$) --- weakly-informative priors or pooling;
(iii) \emph{nonstationarity}: a regime change makes train-window
standardization and as-of joins load onto the intercept; held-out
covariates outside the training support explode $e^{\beta^\top X}$ ---
winsorize predictions to the training envelope (implemented in the stable
fitter);
(iv) \emph{publication lag leakage} if covariates are joined on reference
date rather than release date.

% =======================================================================
\section{Step 2: multi-state models with constrained stage transitions}
\label{sec:step2}

\subsection{The state space and the marked process}
Let the stage lattice be $\mathcal S=\{$Seed$,A,B,C,D,E,\dots\}
\cup\{\partial_{\mathrm{IPO}},\partial_{\mathrm{M\&A}},
\partial_{\dagger}\}$ with the partial order of progression and the
absorbing boundary $\partial$. Firm $i$ occupies state $S_i(t)\in\mathcal S$;
each funding event is a \emph{transition mark}: a jump $k\to j$ with the
admissibility constraint
\begin{equation}
j \succeq k:\qquad
\text{from } C:\; \{C, D, E,\dots\} \text{ allowed (repeats, skips)};\quad
A, B \text{ forbidden}.
\label{eq:admissible}
\end{equation}
The pair $(N_i,\{S_i\})$ is a marked counting process; equivalently a
\emph{progressive multi-state model with self-transitions}
\citep{putter2007,abgk1993}. Absorbing exits make graduation and death part
of the state space instead of censoring nuisances --- the structural cure
for the silent-exit bias of \S\ref{sec:censoring} whenever exit is at least
partially observed.

\subsection{Transition intensities, competing risks, and factorization}
Give every admissible transition its own Aalen multiplicative intensity
\begin{equation}
\lambda_{k\to j,i}(t)\;=\;Y_{ik}(t)\;
\exp\!\big(\alpha_{k\to j}+\beta_{k\to j}^\top X(t)
+\delta_{k\to j}^\top Z_{i,t}\big),
\qquad Y_{ik}(t)=\one\{S_i(t-)=k\},
\label{eq:msm}
\end{equation}
From state $k$ the destinations $\{j\succeq k\}\cup\partial$ \emph{compete}.
Two classical facts organize everything:

\begin{proposition}[Likelihood factorization, ABGK]
If the parameter sets of distinct transitions are variation-independent,
the full likelihood \eqref{eq:pp_lik} of the marked process factorizes over
transitions; each factor is the likelihood of a univariate counting process
in which all other transitions act as censoring. Estimation therefore
decomposes into independent concave fits per transition.
\end{proposition}

\begin{proposition}[Cause-specific vs subdistribution]
The cause-specific hazards in \eqref{eq:msm} are estimable by treating
competing transitions as censoring \emph{regardless of dependence} among
causes \citep{prentice1978}. But probabilities of eventual destinations
(cumulative incidence) require \emph{all} competing hazards, combined by
the Aalen--Johansen product-integral \citep{aalenjohansen1978}; a
Fine--Gray subdistribution model \citep{finegray1999} is the alternative
when the incidence scale is the deliverable. For our ranking use we need
hazards (the mark factor conditions on an event), so cause-specific is the
right scale; for lifetime stage-mix forecasts, Aalen--Johansen.
\end{proposition}

\subsection{Clocks: Markov, semi-Markov, and both}
A modeling fork with real consequences. In \emph{calendar time} (clock
forward), \eqref{eq:msm} is Markov given covariates and macro effects enter
naturally. In \emph{gap time} (clock reset at each round; semi-Markov), the
baseline hazard depends on the time since the last transition ---
empirically the dominant timescale (raise cadence), and exactly what the
renewal literature and the gap-time model of
Prentice--Williams--Peterson \citep{pwp1981} formalize. The pragmatic
resolution, standard in recurrent-event survival, is calendar-time
intensity with gap-time \emph{covariates} --- which is precisely the bridge
to step 3: the semi-Markov content of the process is carried by a
predictable covariate rather than by the clock.

\subsection{Censoring in the multi-state model}
Administrative and observed-exit censoring are handled by
$Y_{ik}$ as in \S\ref{sec:censoring}. Two specifics deserve care:
(i) \emph{left truncation}: a firm first observed at stage $C$ enters the
risk sets of $C\to\cdot$ transitions at its entry time, contributing no
earlier exposure --- delayed entry, standard but mandatory
\citep{kalbfleisch2002}; its unobserved earlier path is irrelevant for
cause-specific hazards by the factorization, a quiet but important
robustness of step 2;
(ii) \emph{silent exits}: still the weak point. Within the multi-state
frame, the honest device is a partially observed absorbing transition: a
latent $k\to\partial_\dagger$ with a cure/mover--stayer structure
\citep{farewell1982}, identified from long silences only under a parametric
assumption --- state it, fit it, and report sensitivity, or bound as in
\S\ref{sec:censoring}.

\subsection{Assumptions}
\begin{assumption}[Admissibility]\label{as:msm_dag}
Transitions violating \eqref{eq:admissible} have intensity zero. (Encoded,
not estimated; data violations are label noise to be measured, not signal.)
\end{assumption}
\begin{assumption}[Variation-independent parameters]\label{as:msm_sep}
No parameter sharing across transitions unless imposed deliberately as
shrinkage (below).
\end{assumption}
\begin{assumption}[Independent censoring given covariates, per
transition]\label{as:msm_cens}
For each transition, \eqref{eq:indep_censoring} holds conditional on
$(X,Z)$.
\end{assumption}

\subsection{What the data must show, and where calibration fails}
\emph{To see in data:} an empirical transition matrix respecting
\eqref{eq:admissible} up to a small, quantified label-noise mass in the
forbidden cells; per-transition event counts sufficient for the cells you
keep unpooled; stage-mix forecasts (Aalen--Johansen) calibrated by cohort.

\emph{Failure modes:}
(i) \emph{sparse cells}: Seed$\to$E has a handful of events; unpenalized
per-transition fits explode. Hierarchical shrinkage --- e.g.\
$\delta_{k\to j}=\delta_0+\Delta_{k\to j}$ with a prior on $\Delta$ ---
or pooling across destinations within origin is mandatory;
(ii) \emph{ties}: weekly data put many transitions at the same timestamp;
use Efron's approximation over Breslow when partial likelihoods are used
\citep{efron1977};
(iii) \emph{stage mislabels} put firms in wrong risk sets --- the forbidden
cells of the transition matrix estimate the mislabel rate, an
error-in-state problem analogous to misclassification in hidden Markov
models;
(iv) \emph{immortal-time bias} if stage is backfilled from later
information --- stage must be as-of, like every covariate.

% =======================================================================
\section{Step 3: history as predictable covariates
(``Hawkes without Hawkes'')}
\label{sec:step3}

\subsection{The model and why it is licensed}
Augment \eqref{eq:msm} with functionals of the past:
\begin{equation}
\lambda_{k\to j,i}(t)=Y_{ik}(t)\exp\Big(
\alpha_{k\to j}+\beta^\top X(t)
+\phi^{\top} H_{s(i)}(t)
-\psi\, g\big(t-T_i^{\mathrm{last}}\big)\Big),
\label{eq:step3}
\end{equation}
where $H_s(t)=\int_0^{t-} \kappa_h(t-u)\dd N_s(u)$ is a
\emph{fixed-kernel} convolution of same-sector activity (e.g.\ exponential
kernels at half-lives 30/90/180 days, coefficients $\phi$ free) and $g$ a
monotone gap transform (log-gap, or binned dummies to let the data draw the
inhibition curve).

The decisive technical point: $H_s$ and $g(t-T^{\mathrm{last}}_i)$ are
\emph{predictable} processes --- measurable with respect to $\cF_{t-}$.
The entire likelihood theory of \S\ref{sec:prelim} is agnostic to whether a
predictable covariate is exogenous or built from the process's own past;
this is the content of the Andersen--Gill formulation of recurrent events
\citep{andersengill1982}. Concavity survives because the parameters still
enter affinely. Step 3 therefore costs \emph{nothing} statistically
relative to step 2, while capturing the two interaction effects that
motivate self-excitation:
sector momentum with sign-free coefficients, and self-inhibition --- which
a nonnegative Hawkes kernel structurally cannot express.

\begin{remark}[Relation to Hawkes]
Equation \eqref{eq:step3} is a multiplicative (log-link) analogue of a
Hawkes process with \emph{known} kernel shape and \emph{estimated}
amplitude. Two genuine differences from step 4: (a) additivity --- a linear
Hawkes intensity is $\mu+\sum\text{kernels}$, supporting the cluster
decomposition of \S\ref{sec:step4}, whereas \eqref{eq:step3} multiplies;
for small exposures $e^{\phi H}\approx 1+\phi H$ and the two coincide
locally; (b) the kernel timescale is fixed here and free there. Fixing it
is not a concession: it is the response to the $\kappa$--$\theta$
identifiability ridge (\code{MATH\_NOTES.md} M3) --- estimate what the data
can identify (amplitudes), profile over a small bank of interpretable
half-lives, and leave joint decay estimation to the regime where it is
identified (\S\ref{sec:step4}).
\end{remark}

\subsection{The deepest calibration problem on the ladder:
heterogeneity vs state dependence}
A positive $\hat\phi$ has two observationally similar sources: true
dynamics (past events genuinely raise the hazard) and \emph{unobserved
heterogeneity} (good sectors/firms have persistently high hazards, hence
clustered events). This is Heckman's classical distinction
\citep{heckman1981}; in duration analysis, frailty and duration dependence
are not separately identified from single spells without covariates
\citep{elbersridder1982,lancaster1990}. Two structural facts rescue us:
\begin{enumerate}[leftmargin=1.6em]
\item \emph{Recurrent events.} Firms and sectors have many spells. With
multi-spell data a shared frailty $u_i$ (or $u_s$) multiplying the
intensity --- gamma frailty being the tractable canonical choice
\citep{vaupel1979,nielsen1992} --- is identified alongside $\phi$: the
frailty explains \emph{permanent} differences in level, the history terms
explain \emph{transient} post-event elevation. Fit \eqref{eq:step3} with
and without frailty; the drop in $\hat\phi$ measures how much
``momentum'' was heterogeneity.
\item \emph{Common factors.} A serially correlated \emph{market-wide}
latent factor is the worst confounder because it is shared across sectors:
it loads onto every $H_s$. The guard is observable proxies (aggregate
market activity as a covariate) plus the diagnostic: if adding the market
aggregate collapses $\hat\phi$, the ``sector momentum'' was a common
factor. Our synthetic rehearsals measured exactly this mechanism inflating
fitted excitation by a factor $\sim2.4$.
\end{enumerate}

\subsection{Assumptions, data requirements, failure modes}
\begin{assumption}[Kernel bank]\label{as:s3_kernel}
The true history dependence is spanned (approximately) by the chosen
fixed-decay kernels; adequacy is checked by the binned-gap refit
(M6) and by residual clustering diagnostics.
\end{assumption}
\begin{assumption}[Frailty orthogonality]\label{as:s3_frailty}
Unobserved heterogeneity is time-invariant (frailty), not itself dynamic;
dynamic unobservables beyond the included factors are attributed to the
history terms and reported as such.
\end{assumption}

\emph{To see in data:} PIT/dispersion clean-up relative to step 2;
stability of $\hat\phi$ across half-life choices (a coefficient that
migrates wildly across the bank means the timescale is wrong); frailty
variance materially $>0$ (expect it); $\hat\psi>0$ with the binned curve
monotone.

\emph{Failure modes:}
(i) frailty--momentum confounding as above --- always fit the frailty
variant;
(ii) \emph{simultaneity} in $H_s$: the firm's own event is in the sector
count --- exclude own events from $H_{s(i)}$ or the coefficient absorbs a
mechanical reflection;
(iii) collinearity between $H_s$ at adjacent half-lives (correlations
$>0.9$) --- orthogonalize or keep two well-separated scales;
(iv) boundary spells: gaps censored at entry (left truncation) must use the
entry-adjusted gap, else early-sample hazards are biased.

% =======================================================================
\section{Step 4: self-exciting processes and the cascade representation}
\label{sec:step4}

\subsection{Linear Hawkes and the cluster/cascade picture}
The linear multivariate Hawkes process \citep{hawkes1971} has intensity
\begin{equation}
\lambda_m(t)\;=\;\mu_m(t)\;+\;\sum_{r}\int_0^{t-}
\varphi_{mr}(t-u)\,\dd N_r(u),
\qquad \varphi_{mr}\ge0 ,
\label{eq:hawkes}
\end{equation}
with branching matrix $K_{mr}=\int_0^\infty\varphi_{mr}$ and stationarity
iff $\rho(K)<1$ \citep{hawkes1971,bremaud1996}. The equivalent
\emph{cluster representation} \citep{hawkesoakes1974} is the conceptual
key: immigrants arrive as a Poisson process with rate $\mu$, and each event
(immigrant or not) independently spawns offspring as an inhomogeneous
Poisson process with rate $\varphi(\cdot)$ --- a branching cascade of
Poisson processes. This is precisely the modeling frame of
\citet{simmajordan2010}, \emph{Modeling Events with Cascades of Poisson
Processes}: events are attributed to a background process or to a parent
event, the latent parentage is marginalised or imputed, and estimation
proceeds by EM over the branching structure --- the same augmentation used
in seismology \citep{veen2008} and equivalent to the complete-data
likelihood being that of independent Poisson cascades. Three consequences
we use:
\begin{enumerate}[leftmargin=1.6em]
\item \emph{EM estimation.} The E-step computes parentage probabilities
$p_{kl}\propto\varphi(t_k-t_l)$ vs $p_{k0}\propto\mu(t_k)$; the M-step
fits Poisson rates on the attributed streams. Monotone, stable, and it
scales; it is the alternative of choice to direct maximization of
\eqref{eq:pp_lik} \citep{ogata1978,ozaki1979} when kernels are flexible
\citep{lewis2011}.
\item \emph{The immigrant/offspring decomposition is meaningful.} The
background stream is covariate-driven; the offspring stream carries the
endogenous dynamics. In funding data we can \emph{partially observe} this
latent attribution: first financings are immigrants by definition. Step 5
exploits exactly this.
\item \emph{Second-order non-identifiability.} A doubly-stochastic
(Cox) process with serially correlated random intensity can match the
second-order statistics (Bartlett spectrum) of a Hawkes process
\citep{daleyvj2003}; counts alone therefore cannot distinguish contagion
from common factors --- the process-level restatement of the
heterogeneity problem of \S\ref{sec:step3}, and the reason every fitted
branching ratio is reported as an upper bound.
\end{enumerate}

\subsection{Estimation regimes and interval censoring}
With exact day-resolution times: direct MLE of \eqref{eq:pp_lik} (concave
in $(\mu,\alpha)$ for fixed decays; profile over decay banks) or EM as
above; goodness of fit by the time-rescaling theorem --- transformed waits
$\Lambda(t_k)-\Lambda(t_{k-1})$ must be i.i.d.\ Exp(1)
\citep{brown2002}. With counts only: the interval-censored route through
the mean-behavior process of \citet{rizoiu2022} (our MBPP machinery), which
identifies $\kappa$ well and the decay poorly --- the Fisher-information
collapse of \code{MATH\_NOTES.md} M3 quantifies the loss as bucket width
grows (safe at $\le$ one kernel half-life, destroyed by $\ge4$).
Discrete-time log-link approximations (Poisson autoregressions on lagged
counts) are \emph{not} Hawkes: their stability theory is of
Fokianos--Tj{\o}stheim type \citep{fokianos2011} and raw-count exponential
feedback is explosive even with small coefficients --- the ``Mechanism B''
pathology we measured; simulation guards are mandatory.

\subsection{Assumptions, data requirements, failure modes}
\begin{assumption}[Subcriticality]\label{as:h_stab}
$\rho(K)<1$, imposed in estimation (spectral/row-sum constraint) and
verified before any forward simulation.
\end{assumption}
\begin{assumption}[Kernel family]\label{as:h_kernel}
$\varphi$ exponential or a small sum of exponentials (Markovian state-space
reduction exists); richer shapes only with EM + volume.
\end{assumption}
\begin{assumption}[No unmodeled common factor]\label{as:h_factor}
Interpreting $\hat K$ structurally requires the common-factor channel
closed (observable proxies included); otherwise $\hat K$ is predictive,
and an upper bound.
\end{assumption}

\emph{To see in data:} day resolution; enough events per modeled node
(the ridge scales as \code{MATH\_NOTES.md} M3 --- thousands of events per
free $(\kappa,\theta)$ pair); time-rescaling KS passing; stability margin
$\rho(\hat K)$ comfortably below 1.

\emph{Failure modes:}
(i) the $\kappa$--$\theta$ ridge (fix or bank the decays);
(ii) common-factor inflation of $\hat K$ (see
Assumption~\ref{as:h_factor});
(iii) boundary solutions: row sums pinned at a stability constraint mean
the likelihood wants more feedback than stationarity allows --- a
regime-change or latent-factor signal, not contagion;
(iv) edge effects: cascades truncated at $t=0$ bias early-sample kernels
(condition on a burn-in window);
(v) multi-timescale confusion: daily reporting spikes (Monday batching of
deal announcements) masquerade as fast kernels --- deseasonalize the
within-week pattern first.

% =======================================================================
\section{Step 5: the outside universe --- a two-step nested structure}
\label{sec:step5}

\subsection{The mark factor with an outside option}
Conditional on an event in sector $s$ at $t$, the winner is either a member
of the observed risk set $\cR_{s,t}$ or ``outside'' (a first-time entrant,
or an untracked incumbent). Write $O$ for the outside alternative. The
natural structure is a \emph{nested logit} \citep{mcfadden1974,
mcfadden1978}: an upper choice between nests $\{O\}$ and $\cR_{s,t}$, then
a within-nest choice among incumbents,
\begin{align}
\Prob(\text{outside}\mid s,t)
&=\sigma\!\big(b_0+\gamma^\top g_{s,t}\big),
\label{eq:step5_upper}\\[2pt]
\Prob(i\mid \text{inside}, s,t)
&=\frac{\exp\big(q_{i,t}\big)}
{\sum_{j\in\cR_{s,t}}\exp\big(q_{j,t}\big)},
\label{eq:step5_lower}
\end{align}
with $g_{s,t}$ \emph{global/context} covariates only (sector activity,
historical first-financing share, macro, pool size) --- necessarily so,
because firm covariates do not exist for the outside branch. The joint
log-likelihood is the exact sum of the two stages; with a free upper-stage
model this is the nested logit at dissimilarity parameter 1, and the
within-nest estimates are consistent regardless of the upper stage (the
two-step estimator of \citealp{mcfadden1978}), which also licenses
\emph{sampling of alternatives} in the big denominators of
\eqref{eq:step5_lower} --- uniform sampling with the winner retained leaves
the conditional MLE consistent (\code{MATH\_NOTES.md} M4).

\subsection{Three deep readings of the outside option}
\paragraph{(a) Demand-estimation reading.} In differentiated-products demand
\citep{berry1994}, the outside good's share is what identifies market size
and pins utility levels. Identically here: the newcomer share over time is
the observable that identifies \emph{coverage} --- a drifting outside share
under stable context covariates is the diagnostic that the universe, not
the market, changed. And the converse identifiability warning holds
(\code{MATH\_NOTES.md} M7): a pool-size term in the outside utility is not
separately identified from the intercept unless pool size varies over time.
\paragraph{(b) Truncation reading.} Entrants are left-truncated at first
funding: their covariates are unobservable pre-entry, so any within-nest
score for them must integrate over the covariate distribution rather than
impute. Under a Gaussian entrant law $z\sim N(\mu_s,\Sigma_s)$ the
integrated utility is the log-MGF $w^\top\mu_s+\tfrac12 w^\top\Sigma_s w$
--- exact, estimable from observed first deals, and a measured
decomposition device (it attributes the fitted outside constant into
quality vs coverage; empirically quality was the minor share).
\paragraph{(c) Cascade reading.} In the cluster representation of
\S\ref{sec:step4}, the outside stream \emph{is} the immigrant process:
covariate-driven background arrivals, versus offspring arrivals from within
the observed system. Simma--Jordan's latent parentage becomes partially
observed in funding data --- a first financing is an immigrant with
certainty --- so the two-step is not an engineering workaround; it is the
observable projection of the branching decomposition, estimated without EM
precisely because the data label the immigrants.

\subsection{Coarsening, CAR, and the conflation}
The outside label conflates two populations (true entrants; untracked
incumbents). Formally the mark is \emph{coarsened}: we observe the coarse
label $O$ instead of the fine identity. The two-step likelihood is valid
under Coarsening At Random \citep{heitjan1991} --- the probability of being
untracked must not depend on the fine identity beyond what $g_{s,t}$
carries. CAR here is not innocuous (coverage correlates with sector and
time), which is exactly why the upper stage carries context covariates:
they move the coarsening mechanism into the conditioning set. What is lost
under conflation is \emph{composition}, not fit --- separating the classes
with oracle labels moved the joint likelihood negligibly in rehearsal ---
so the discipline is linguistic as much as statistical: report ``outside
probability,'' never ``entrant probability.''

\subsection{Assumptions, data requirements, failure modes}
\begin{assumption}[Upper-stage sufficiency]\label{as:s5_car}
Given $g_{s,t}$, the outside/inside indicator is independent of the fine
identity of the winner (CAR given context).
\end{assumption}
\begin{assumption}[Within-nest correctness]\label{as:s5_nest}
Conditional on inside, the risk set is correct --- all of
\S\ref{sec:censoring} item 3 applies to $\cR_{s,t}$, and its failure modes
dominate any nesting refinement.
\end{assumption}

\emph{To see in data:} quarterly reliability of the predicted outside share
(a few points of calibration error at most); stability of within-nest
coefficients when the upper stage is re-specified (they should not move ---
a factorization property that doubles as a specification test); entrant
first-deal feature distributions stable enough across the train window for
the log-MGF device.

\emph{Failure modes:}
(i) coverage drift misread as market change (watch the context-adjusted
outside share);
(ii) upper-stage separation in sector-quarters with no outside wins;
(iii) IIA violations \emph{within} the incumbent nest (a red-flag test:
estimated dissimilarity far from 1 when relaxed, or systematic
substitution-pattern residuals) --- the escalation is a richer nest
structure, not abandonment of the frame;
(iv) the pool-size/intercept non-identifiability (M7) --- never interpret
$b_0$ alone.

% =======================================================================
\section{Synthesis: the assumptions ladder and the calibration checklist}
\label{sec:synthesis}

Each rung strictly weakens the independence structure of the last:
\[
\underbrace{\text{Poisson}(X)}_{\text{no history}}
\;\subset\;
\underbrace{\text{multi-state}}_{\text{history via state}}
\;\subset\;
\underbrace{\text{+ predictable history covariates}}_{\text{fixed-kernel dynamics}}
\;\subset\;
\underbrace{\text{Hawkes}}_{\text{free-kernel dynamics}}
\]
with step 5 orthogonal (a statement about the \emph{mark} space, attachable
to any rung). The corresponding falsification chain runs in the same order:
dispersion/PIT reject rung 1; forbidden-cell mass and stage-mix
calibration test rung 2; residual clustering after fixed kernels is the
only license for rung 4; outside-share reliability tests rung 5. Every
assumption above is stated so that a specific diagnostic can break it ---
that, more than any single model, is the methodological commitment: the
ladder is climbable in either direction, and the data decide the floor.

% =======================================================================
\begin{thebibliography}{99}\small

\bibitem[Aalen(1978)]{aalen1978}
Aalen, O.~O. (1978). Nonparametric inference for a family of counting
processes. \emph{Annals of Statistics} 6(4), 701--726.

\bibitem[Aalen \& Johansen(1978)]{aalenjohansen1978}
Aalen, O.~O., \& Johansen, S. (1978). An empirical transition matrix for
non-homogeneous Markov chains based on censored observations.
\emph{Scandinavian Journal of Statistics} 5, 141--150.

\bibitem[Andersen \& Gill(1982)]{andersengill1982}
Andersen, P.~K., \& Gill, R.~D. (1982). Cox's regression model for counting
processes: a large sample study. \emph{Annals of Statistics} 10(4),
1100--1120.

\bibitem[Andersen et al.(1993)]{abgk1993}
Andersen, P.~K., Borgan, {\O}., Gill, R.~D., \& Keiding, N. (1993).
\emph{Statistical Models Based on Counting Processes}. Springer.

\bibitem[Berry(1994)]{berry1994}
Berry, S.~T. (1994). Estimating discrete-choice models of product
differentiation. \emph{RAND Journal of Economics} 25(2), 242--262.

\bibitem[Br\'emaud \& Massouli\'e(1996)]{bremaud1996}
Br\'emaud, P., \& Massouli\'e, L. (1996). Stability of nonlinear Hawkes
processes. \emph{Annals of Probability} 24(3), 1563--1588.

\bibitem[Brown et al.(2002)]{brown2002}
Brown, E.~N., Barbieri, R., Ventura, V., Kass, R.~E., \& Frank, L.~M.
(2002). The time-rescaling theorem and its application to neural spike
train data analysis. \emph{Neural Computation} 14(2), 325--346.

\bibitem[Cox(1972)]{cox1972}
Cox, D.~R. (1972). Regression models and life-tables (with discussion).
\emph{Journal of the Royal Statistical Society B} 34(2), 187--220.

\bibitem[Daley \& Vere-Jones(2003)]{daleyvj2003}
Daley, D.~J., \& Vere-Jones, D. (2003). \emph{An Introduction to the Theory
of Point Processes}, Vol.~I, 2nd ed. Springer.

\bibitem[Efron(1977)]{efron1977}
Efron, B. (1977). The efficiency of Cox's likelihood function for censored
data. \emph{JASA} 72(359), 557--565.

\bibitem[Elbers \& Ridder(1982)]{elbersridder1982}
Elbers, C., \& Ridder, G. (1982). True and spurious duration dependence:
the identifiability of the proportional hazard model. \emph{Review of
Economic Studies} 49(3), 403--409.

\bibitem[Farewell(1982)]{farewell1982}
Farewell, V.~T. (1982). The use of mixture models for the analysis of
survival data with long-term survivors. \emph{Biometrics} 38(4),
1041--1046.

\bibitem[Fine \& Gray(1999)]{finegray1999}
Fine, J.~P., \& Gray, R.~J. (1999). A proportional hazards model for the
subdistribution of a competing risk. \emph{JASA} 94(446), 496--509.

\bibitem[Fokianos \& Tj{\o}stheim(2011)]{fokianos2011}
Fokianos, K., \& Tj{\o}stheim, D. (2011). Log-linear Poisson
autoregression. \emph{Journal of Multivariate Analysis} 102(3), 563--578.

\bibitem[Hawkes(1971)]{hawkes1971}
Hawkes, A.~G. (1971). Spectra of some self-exciting and mutually exciting
point processes. \emph{Biometrika} 58(1), 83--90.

\bibitem[Hawkes \& Oakes(1974)]{hawkesoakes1974}
Hawkes, A.~G., \& Oakes, D. (1974). A cluster process representation of a
self-exciting process. \emph{Journal of Applied Probability} 11(3),
493--503.

\bibitem[Heckman(1981)]{heckman1981}
Heckman, J.~J. (1981). Heterogeneity and state dependence. In
\emph{Studies in Labor Markets}, University of Chicago Press, 91--140.

\bibitem[Heitjan \& Rubin(1991)]{heitjan1991}
Heitjan, D.~F., \& Rubin, D.~B. (1991). Ignorability and coarse data.
\emph{Annals of Statistics} 19(4), 2244--2253.

\bibitem[Kalbfleisch \& Prentice(2002)]{kalbfleisch2002}
Kalbfleisch, J.~D., \& Prentice, R.~L. (2002). \emph{The Statistical
Analysis of Failure Time Data}, 2nd ed. Wiley.

\bibitem[Lancaster(1990)]{lancaster1990}
Lancaster, T. (1990). \emph{The Econometric Analysis of Transition Data}.
Cambridge University Press.

\bibitem[Lewis \& Mohler(2011)]{lewis2011}
Lewis, E., \& Mohler, G. (2011). A nonparametric EM algorithm for
multiscale Hawkes processes. \emph{Journal of Nonparametric Statistics}
1(1), 1--20.

\bibitem[McFadden(1974)]{mcfadden1974}
McFadden, D. (1974). Conditional logit analysis of qualitative choice
behavior. In \emph{Frontiers in Econometrics}, Academic Press, 105--142.

\bibitem[McFadden(1978)]{mcfadden1978}
McFadden, D. (1978). Modeling the choice of residential location. In
\emph{Spatial Interaction Theory and Planning Models}, North-Holland,
75--96.

\bibitem[Nielsen et al.(1992)]{nielsen1992}
Nielsen, G.~G., Gill, R.~D., Andersen, P.~K., \& S{\o}rensen, T.~I.~A.
(1992). A counting process approach to maximum likelihood estimation in
frailty models. \emph{Scandinavian Journal of Statistics} 19(1), 25--43.

\bibitem[Ogata(1978)]{ogata1978}
Ogata, Y. (1978). The asymptotic behaviour of maximum likelihood
estimators for stationary point processes. \emph{Annals of the Institute
of Statistical Mathematics} 30(1), 243--261.

\bibitem[Ozaki(1979)]{ozaki1979}
Ozaki, T. (1979). Maximum likelihood estimation of Hawkes' self-exciting
point processes. \emph{Annals of the Institute of Statistical
Mathematics} 31(1), 145--155.

\bibitem[Prentice et al.(1978)]{prentice1978}
Prentice, R.~L., Kalbfleisch, J.~D., Peterson, A.~V., Flournoy, N.,
Farewell, V.~T., \& Breslow, N.~E. (1978). The analysis of failure times
in the presence of competing risks. \emph{Biometrics} 34(4), 541--554.

\bibitem[Prentice et al.(1981)]{pwp1981}
Prentice, R.~L., Williams, B.~J., \& Peterson, A.~V. (1981). On the
regression analysis of multivariate failure time data. \emph{Biometrika}
68(2), 373--379.

\bibitem[Putter et al.(2007)]{putter2007}
Putter, H., Fiocco, M., \& Geskus, R.~B. (2007). Tutorial in
biostatistics: competing risks and multi-state models. \emph{Statistics in
Medicine} 26(11), 2389--2430.

\bibitem[Rizoiu et al.(2022)]{rizoiu2022}
Rizoiu, M.-A., Soen, A., Li, S., Calderon, P., Dong, L., Menon, A.~K., \&
Xie, L. (2022). Interval-censored Hawkes processes. \emph{Journal of
Machine Learning Research} 23(338), 1--84.

\bibitem[Simma \& Jordan(2010)]{simmajordan2010}
Simma, A., \& Jordan, M.~I. (2010). Modeling events with cascades of
Poisson processes. In \emph{Proceedings of the 26th Conference on
Uncertainty in Artificial Intelligence (UAI)}, 546--555.

\bibitem[Vaupel et al.(1979)]{vaupel1979}
Vaupel, J.~W., Manton, K.~G., \& Stallard, E. (1979). The impact of
heterogeneity in individual frailty on the dynamics of mortality.
\emph{Demography} 16(3), 439--454.

\bibitem[Veen \& Schoenberg(2008)]{veen2008}
Veen, A., \& Schoenberg, F.~P. (2008). Estimation of space--time branching
process models in seismology using an EM-type algorithm. \emph{JASA}
103(482), 614--624.

\bibitem[Wei et al.(1989)]{wlw1989}
Wei, L.~J., Lin, D.~Y., \& Weissfeld, L. (1989). Regression analysis of
multivariate incomplete failure time data by modeling marginal
distributions. \emph{JASA} 84(408), 1065--1073.

\end{thebibliography}

\end{document}
